\section{Formally Real: Square Roots}

\subsection{Square Root Predicate}

\begin{definition}[IsPositiveSqrt]
Let $a, b \in J$. We say $b$ is a \textbf{positive square root} of $a$ if
\[ \jsq{b} = a \quad \land \quad b \in \mathsf{PositiveElement}(J). \]
\end{definition}

\begin{definition}[HasPositiveSqrt]
An element $a \in J$ \textbf{has a positive square root} if
\[ \exists\, b \in J \text{ such that } \mathsf{IsPositiveSqrt}(a, b). \]
\end{definition}

\subsection{Basic Properties}

\begin{lemma}[isPositiveSqrt\_def]
For $a, b \in J$:
\[ \mathsf{IsPositiveSqrt}(a, b) \iff \jsq{b} = a \land \mathsf{PositiveElement}(b). \]
\end{lemma}

\begin{theorem}[PositiveElement.of\_hasPositiveSqrt]
If $a \in J$ has a positive square root, then $a$ is a positive element.
\end{theorem}

\begin{lemma}[isPositiveSqrt\_zero]
Zero has zero as its positive square root:
\[ \mathsf{IsPositiveSqrt}(0, 0). \]
\end{lemma}

\begin{lemma}[isPositiveSqrt\_jone]
If $\jone$ is a positive element, then $\jone$ is its own positive square root:
\[ \mathsf{PositiveElement}(\jone) \implies \mathsf{IsPositiveSqrt}(\jone, \jone). \]
\end{lemma}

\subsection{Uniqueness}

\begin{theorem}[isPositiveSqrt\_unique]
Let $J$ be a formally real Jordan algebra. Let $a, b, c \in J$ with
\begin{align*}
\mathsf{IsPositiveSqrt}(a, b), \quad &\mathsf{IsPositiveSqrt}(a, c), \\
b \jmul c = c \jmul b, \quad &\jsq{b} \jmul c = c \jmul \jsq{b}.
\end{align*}
Then $b = c$.
\end{theorem}

\subsection{Existence}

\begin{theorem}[HasPositiveSqrt.of\_positiveElement]
Let $a \in J$ be a positive element. Then $a$ has a positive square root:
\[ \mathsf{PositiveElement}(a) \implies \mathsf{HasPositiveSqrt}(a). \]
\end{theorem}
